\documentclass[a4paper]{article}

\usepackage{graphicx}
\usepackage{amsmath,amssymb}
\usepackage{minted}
\usepackage{multirow}
\usepackage{float}
\usepackage{algorithm}
\usepackage{algpseudocode}
\usepackage{todonotes}
\usepackage{forest}
\usepackage{xstring}
\usepackage{xcolor}
\usepackage[most]{tcolorbox}
\usepackage{xparse}
\usepackage{natbib}

\usetikzlibrary{graphs,graphdrawing}
\usegdlibrary{layered}

\title{Mind Mirror Release History: Native Source Layout Crash}
\date{2026-09-04}
\author{Mohammad Rahmani}

\begin{document}
   \maketitle

   \section{Bug}
   Mind Mirror exited with native process code \texttt{-11} immediately after
   startup. The control-panel log contained only the launch header and no Python
   traceback because the process terminated inside Qt before Python exception
   handling could report an error.

   The regression was introduced by installing a Python subclass of
   \texttt{QPlainTextDocumentLayout}. Its overridden block-layout path modified
   \texttt{QTextLayout} while Qt was already performing native plain-text layout,
   creating an unsafe re-entrant layout path.

   \section{Resolution}
   \begin{itemize}
      \item Removed the custom Python \texttt{QPlainTextDocumentLayout} from the
      active Source editor architecture.
      \item Restored Qt's native \texttt{QTextEdit}/\texttt{QTextDocumentLayout}
      path for Source presentation. This again honors block direction, alignment,
      proportional line height, and the block top margins used by inline Source
      graphic previews.
      \item Persian and Arabic Source prose therefore uses native right-to-left
      block direction and right alignment, while LaTeX control lines remain
      left-to-right according to the existing direction resolver.
      \item Large LaTeX sources continue to use the chunked canonical
      \texttt{LatexSourceBuffer}. Once the configured large-document threshold is
      reached, only an approximately \texttt{40,000}-character Source window is
      materialized in \texttt{QTextEdit}.
      \item The incremental \texttt{QTextDocument.contentsChange} path is retained:
      normal typing updates only the edited source range in the chunk buffer and
      does not serialize the complete Source window on every key press.
      \item Large mode omits decorative hanging-indent calculation, indent-guide
      painting, and soft-wrap-marker painting while preserving native RTL and
      figure reservation geometry.
   \end{itemize}

   \section{Regression coverage}
   Source-level regression tests verify that the public Source editor no longer
   installs a custom plain-text document layout, that the native rich editor
   retains RTL block formatting, that Source figures reserve block top margins,
   and that the large Source editor keeps the incremental chunk-edit path while
   materializing a bounded rich-text Source window.

   \bibliographystyle{plainnat}
   \bibliography{/home/donkarlo/Dropbox/repo/nd_shared_research/bibliography/bibliography.bib}
\end{document}
